% This file demonstrates how to use the IEEEConf LaTeX2e macro package,
% to prepare a manuscript for proceedings on CD of the conference
% FedCSIS
%
\documentclass[conference]{IEEEtran}
\IEEEoverridecommandlockouts

% This package serves to balance the column lengths on the last page of the document.
\usepackage{pbalance}

%% to enable \thank command
\IEEEoverridecommandlockouts 
%% The usage of the following packages is recommended
%% to insert graphics
\usepackage[dvips]{graphicx}
% to typeset algorithms
\usepackage{algorithmic}
\usepackage{algorithm}
% to typeset code fragments
\usepackage{minted}
% to make an accent \k be available
\usepackage[OT4,T1]{fontenc}
% provides various features to facilitate writing math formulas and to improve the typographical quality of their output.
\usepackage[cmex10]{amsmath}
\interdisplaylinepenalty=2500
% por urls typesetting and breaking
\usepackage{url}
% for vertical merging table cells
\usepackage{multirow}

% define environments for remarks and examples
\newtheorem{remark}{Remark}[section]
\newtheorem{example}[remark]{Example}

%
%
\title{Federated Learning for Histopathological Image Segmentation: A Privacy-Preserving Study Using the EBHI-Seg Dataset}
%
%
\author{
\IEEEauthorblockN{Mohammadreza Azimi}
\IEEEauthorblockA{0000-0002-9036-2122\\
Riga Stradins University\\
in Riga\\
Ratsupites iela 5 Riga, LV-1067, Latvija\\
Email: Mohammadreza.Azimi@rsu.lv}
\and
\IEEEauthorblockN{Henrik Gabrielyan}
\IEEEauthorblockA{0000-0002-9036-2122\\
Riga Stradins University\\
in Riga\\
Ratsupites iela 5 Riga, LV-1067, Latvija\\
Email: Henrik.Gabrielyan@rsu.lv}
}

% conference papers do not typically use \thanks and this command
% is locked out in conference mode. If really needed, such as for
% the acknowledgment of grants, issue a \IEEEoverridecommandlockouts
% after \documentclass

% for over three affiliations, or if they all won't fit within the width
% of the page, use this alternative format:
% 
%\author{\IEEEauthorblockN{Michael Shell\IEEEauthorrefmark{1},
%Homer Simpson\IEEEauthorrefmark{2},
%James Kirk\IEEEauthorrefmark{3}, 
%Montgomery Scott\IEEEauthorrefmark{3} and
%Eldon Tyrell\IEEEauthorrefmark{4}}
%\IEEEauthorblockA{\IEEEauthorrefmark{1}School of Electrical and Computer Engineering\\
%Georgia Institute of Technology,
%Atlanta, Georgia 30332--0250\\ Email: see http://www.michaelshell.org/contact.html}
%\IEEEauthorblockA{\IEEEauthorrefmark{2}Twentieth Century Fox, Springfield, USA\\
%Email: homer@thesimpsons.com}
%\IEEEauthorblockA{\IEEEauthorrefmark{3}Starfleet Academy, San Francisco, California 96678-2391\\
%Telephone: (800) 555--1212, Fax: (888) 555--1212}
%\IEEEauthorblockA{\IEEEauthorrefmark{4}Tyrell Inc., 123 Replicant Street, Los Angeles, California 90210--4321}}





\begin{document}
\maketitle              % typeset the title of the contribution

\begin{abstract}
The increasing adoption of artificial intelligence in digital pathology is often constrained by limited data sharing due to privacy and regulatory concerns. Federated learning (FL) has emerged as a promising paradigm that enables collaborative model training across institutions without centralizing sensitive data. In this study, we investigate the feasibility and performance of federated learning for histopathological image segmentation using the EBHI-Seg dataset.

We simulate a multi-institutional setting by partitioning the dataset into multiple decentralized clients, reflecting realistic clinical data distributions under both independent and non-identically distributed (non-IID) conditions. A convolutional neural network based on the U-Net architecture is employed as the baseline segmentation model and trained using a federated optimization strategy. Model performance is evaluated using standard metrics, including Dice similarity coefficient and Intersection over Union (IoU), and compared against a centrally trained counterpart.

Our results demonstrate that federated learning achieves competitive segmentation performance while preserving data locality, although performance degradation is observed under highly heterogeneous data distributions. Further analysis highlights the impact of client variability, data imbalance, and domain shifts on model convergence and generalization. These findings underscore the potential of federated learning as a viable framework for privacy-preserving collaborative medical image analysis and provide insights into its practical deployment in digital pathology..
\end{abstract}


\section{Introduction}
%
The integration of artificial intelligence (AI) into digital pathology has significantly advanced the field of computational histopathology, enabling automated analysis of whole-slide images for tasks such as tumor detection, grading, and tissue segmentation. In particular, deep learning models, such as convolutional neural networks (CNNs), have demonstrated remarkable performance in medical image segmentation, with architectures like U-Net becoming a standard due to their ability to capture both local and contextual features ~\cite{1}. Despite these advancements, the success of such models is heavily dependent on access to large, diverse, and well-annotated datasets.

However, in medical domains, data sharing across institutions remains a major challenge. Strict privacy regulations, such as General Data Protection Regulation (GDPR), along with ethical and legal constraints, limit the centralization of patient data ~\cite{2}. Additionally, institutional policies and concerns over data ownership further hinder collaborative model development. As a result, many AI models in digital pathology are trained on limited and potentially biased datasets, which can negatively impact their generalizability and robustness across different clinical settings ~\cite{3}.

To address these limitations, Federated Learning (FL) has emerged as a promising distributed learning framework. Initially proposed by Brendan McMahan et al. ~\cite{4}, FL enables multiple institutions (clients) to collaboratively train a shared global model without exchanging raw data. Instead, model updates are computed locally and aggregated centrally, preserving data privacy while leveraging knowledge from decentralized sources. This paradigm is particularly well-suited for healthcare applications, where data sensitivity and heterogeneity are critical concerns.

Recent studies have explored the application of FL in medical imaging, including radiology and pathology, demonstrating its potential to achieve performance comparable to centralized training under certain conditions ~\cite{5,6}. However, several challenges remain, especially in the context of histopathological image segmentation. These include non-identically distributed (non-IID) data across institutions, variations in staining protocols, scanner differences, and class imbalance, all of which can affect model convergence and performance ~\cite{7}. Moreover, the impact of such heterogeneity on segmentation tasks—compared to classification—remains underexplored.

In this work, we investigate the feasibility of federated learning for histopathological image segmentation using the EBHI-Seg dataset. By simulating a multi-institutional environment with varying data distributions, we aim to systematically evaluate the performance of a federated U-Net model relative to a centrally trained baseline. Our study further analyzes the influence of client variability and domain shifts on model performance, providing practical insights into the deployment of FL systems in real-world digital pathology workflows.


\section{Literature Review}

The rapid advancement of deep learning has significantly improved performance in medical image analysis, particularly in tasks such as classification, detection, and segmentation. Convolutional neural networks (CNNs), especially architectures like U-Net, have become the standard for histopathological image segmentation due to their ability to capture both local and global contextual features. However, these approaches typically rely on centralized datasets, which are often difficult to construct in the medical domain due to privacy concerns, regulatory constraints, and institutional data silos.

To address these limitations, Federated Learning (FL) has emerged as a promising distributed learning framework that enables collaborative model training without sharing raw data. Instead, participating institutions train local models and share only model updates, thereby preserving data privacy while still benefiting from multi-institutional knowledge. Recent surveys have highlighted the growing importance of FL in medical imaging, particularly for overcoming data scarcity and enhancing model generalization across diverse clinical environments ~\cite{LL1}.

Several comprehensive reviews have analyzed the application of FL in medical image analysis. For instance, studies published in Diagnostics and Electronics emphasize that FL enables the use of distributed datasets while addressing privacy and security concerns inherent in healthcare systems ~\cite{LL2}. These works also highlight key challenges, including communication overhead, system heterogeneity, and performance degradation under non-identically distributed (non-IID) data conditions. Similarly, research in radiology demonstrates that FL can facilitate multi-center collaborations without requiring data pooling, making it particularly suitable for sensitive imaging modalities such as MRI and CT scans ~\cite{LL3}.

Despite its advantages, federated learning introduces several technical challenges, especially in medical image segmentation. Variability in imaging protocols, annotation standards, and scanner characteristics across institutions leads to significant domain shifts that negatively impact model performance. Recent systematic reviews have shown that segmentation tasks are particularly sensitive to such heterogeneity, as differences in resolution, staining, and labeling quality can reduce model accuracy and convergence stability ~\cite{LL4}. These issues are further exacerbated in histopathology, where staining variability and tissue morphology differ substantially between laboratories.

To mitigate these challenges, several advanced FL methodologies have been proposed. For example, domain generalization techniques such as those introduced in FedDG aim to improve model robustness across unseen institutions by explicitly modeling inter-client distribution shifts. These approaches leverage episodic training strategies to enhance generalization in decentralized environments. Additionally, cross-domain federated learning frameworks have demonstrated the feasibility of training models across heterogeneous datasets while maintaining acceptable segmentation performance, even in multi-organ and multi-modal scenarios ~\cite{LL5}.

Recent experimental studies further validate the effectiveness of FL in medical image segmentation tasks. For instance, a study published in Scientific Reports demonstrated that federated learning combined with attention-based architectures can achieve competitive performance in brain tumor segmentation, highlighting the potential of FL for complex clinical applications ~\cite{LL6}. Moreover, emerging research on large-scale medical imaging models suggests that FL can support continuous model updates and integration of new data sources without violating privacy constraints, making it suitable for real-world deployment in healthcare systems~\cite{LL7}.

Overall, the literature indicates that federated learning represents a transformative approach for medical image analysis, particularly in scenarios where data sharing is restricted. While significant progress has been made, challenges related to data heterogeneity, communication efficiency, and model generalization remain open research problems. Addressing these limitations is essential for enabling robust and scalable FL-based solutions in digital pathology and other medical imaging domains.


\section{Methods and Materials}
In this section, thechosen dataset and the utilized methodology will be explained in details.

\subsection{Methods}

\subsubsection{Data Partitioning Strategy}
To simulate a realistic multi-institutional environment, the EBHI-Seg dataset is partitioned across six virtual clients using two complementary strategies:

\begin{itemize}
\item \emph{Non-IID Partitioning (by-class):} Each client is assigned data from a single tissue class. This strategy reflects scenarios where different institutions specialize in specific pathological conditions. The six tissue classes (Adenocarcinoma, Low-grade IN, Polyp, Serrated adenoma, Normal, and High-grade IN) are distributed one per client, resulting in significant inter-client heterogeneity in class composition and feature distributions.

\item \emph{IID Partitioning (random):} Samples are randomly shuffled and distributed equally across clients, with each receiving approximately 260 training samples. This strategy simulates an idealized homogeneous setting where data distributions are identical across institutions, providing an upper bound for federated learning performance.
\end{itemize}

Data partitioning occurs after the initial train-validation-test split: 70\% training (1,560 samples), 15\% validation (333 samples), and 15\% test (333 samples).

\subsubsection{Segmentation Models}
Three state-of-the-art segmentation architectures are evaluated within the federated learning framework:

\begin{itemize}
\item \emph{UNet:} A custom lightweight encoder-decoder architecture with 32 base channels, yielding approximately 1.9 million parameters. This model serves as the primary baseline due to its widespread adoption in medical image segmentation and computational efficiency suitable for federated settings.

\item \emph{DeepLabV3:} A ResNet-50-based encoder-decoder architecture with atrous spatial pyramid pooling (ASPP) for multi-scale feature extraction, comprising approximately 39 million parameters. This represents a more complex, state-of-the-art approach capable of capturing multi-scale contextual information.

\item \emph{FCN (Fully Convolutional Networks):} A ResNet-50 FCN32s variant with approximately 35 million parameters, providing another competitive baseline for comparison.
\end{itemize}

All models are initialized from scratch without ImageNet pre-training to ensure a fair comparison across centralized and federated settings.

\subsubsection{Federated Learning Algorithms}
Three federated optimization algorithms are compared to assess their robustness to data heterogeneity:

\begin{itemize}
\item \emph{FedAvg:} The foundational federated averaging algorithm that computes a weighted average of client model updates at the server, with client weights proportional to local training set sizes.

\item \emph{FedProx:} An algorithm designed to address extreme data heterogeneity by introducing a proximal term in the local objective function, encouraging clients to maintain consistency with the global model during local training and mitigating divergence caused by non-IID data.

\item \emph{FedOptimizer:} A momentum-based federated optimization algorithm maintaining momentum both at client and server levels, improving convergence speed and stability in heterogeneous settings through accumulated gradient information.
\end{itemize}

\subsubsection{Training Configuration}
All federated learning experiments employ the following hyperparameter configuration:

\begin{itemize}
\item \emph{Federated Settings:} 20 communication rounds, 3 local SGD epochs per round, 6 clients participating in all rounds
\item \emph{Centralized Baseline:} 30 epochs on pooled dataset
\item \emph{Batch Size:} 16 samples
\item \emph{Learning Rate:} 0.001 with weight decay 0.0001
\item \emph{Loss Function:} Weighted BCE-Dice loss combining binary cross-entropy (weight 0.5) and Dice loss (weight 0.5)
\item \emph{Optimizer:} Stochastic gradient descent with momentum (default momentum 0.9)
\item \emph{Device:} NVIDIA T4 GPU with 15 GB VRAM (Google Colab)
\end{itemize}

\subsubsection{Evaluation Metrics}
Model segmentation performance is measured using standard metrics:

\begin{itemize}
\item \emph{Dice Similarity Coefficient (DSC):} Measures spatial overlap between predicted and ground-truth segmentations, computed as $\text{DSC} = \frac{2TP}{2TP + FP + FN}$ where TP, FP, and FN denote true positives, false positives, and false negatives respectively.

\item \emph{Intersection over Union (IoU):} Computes the Jaccard index as the ratio of intersection to union, defined as $\text{IoU} = \frac{TP}{TP + FP + FN}$. This metric is particularly sensitive to false positives.

\item \emph{Training Loss:} BCE-Dice loss value recorded on validation and test sets throughout training.
\end{itemize}

Metrics are recorded at the final communication round (round 20 for federated experiments) and compared against the centralized baseline to quantify performance degradation introduced by federated learning and data heterogeneity.


\subsection{Materials}
 This study utilizes the EBHI-Seg dataset, introduced in the article by Shi et al ~\cite{dataset}. The dataset consists of high-resolution hematoxylin and eosin stained histopathological images acquired from endoscopic biopsy samples of the gastrointestinal tract. It was specifically curated to support the development and evaluation of machine learning models for semantic segmentation tasks in digital pathology.

The EBHI-Seg dataset comprises a collection of digitized whole-slide image patches accompanied by expert-annotated pixel-level segmentation masks. These annotations delineate relevant tissue structures, including epithelial regions, stromal components, and pathological areas such as tumor regions. The images exhibit considerable variability in terms of tissue morphology, staining intensity, and structural composition, reflecting realistic clinical heterogeneity. This variability makes the dataset particularly suitable for evaluating model robustness and generalization in medical image analysis.

All images in the dataset were prepared following standard histopathological procedures, including formalin fixation, paraffin embedding, sectioning, and staining. The digitization process was performed using high-resolution slide scanners, ensuring sufficient detail for cellular-level analysis. Corresponding ground truth annotations were generated by domain experts, enabling supervised learning approaches for precise segmentation tasks. The availability of these high-quality annotations supports the training and validation of deep learning architectures such as U-Net and its variants.

For the purposes of this study, the dataset is treated within a federated learning framework. Although originally centralized, it is partitioned into multiple subsets to simulate distributed data sources, analogous to independent clinical institutions. Each subset represents a local client with its own data distribution, enabling the investigation of federated optimization under both independent and non-identically distributed (non-IID) conditions. This setup facilitates the assessment of collaborative model training without direct data sharing, aligning with privacy-preserving requirements in medical applications.

To further emulate real-world deployment scenarios, controlled heterogeneity is introduced across the simulated clients by varying sample sizes, class distributions, and image characteristics such as staining variations. This allows for a more realistic evaluation of federated learning performance in the presence of domain shifts commonly observed across different healthcare centers.

\begin{minted}{latex}
\usepackage{pbalance}
\end{minted}

\section{Figures and Tables}
To insert a figure you are encouraged to use the  \verb|graphicx| package:
\begin{minted}{latex}
\usepackage[dvips]{graphicx}
\end{minted}
Figures format should be inserted in the \verb|figure| environment, as follows:
\begin{minted}{latex}
\begin{figure}[tbp]
\centering
\includegraphics[width=0.75\hsize]
     {test.pdf}
\caption{Sinusoid}
\end{figure}
\end{minted}
%
\begin{figure}[tbp]
\centering
\includegraphics[width=0.75\hsize]{test.pdf}
\caption{Sinusoid}
\end{figure}
%
Parameters of the command \mintinline{latex}|\includegraphics| are described in the documentation of the \mintinline{latex}|graphicx| package.

Prepare your figures using vector graphics and converting it to standard pdf format. Photo files should be in high-resolution (at least 300 DPI) JPEG files. Consider using of \LaTeX\ packages (TeXDraw, XYpic, TikZ, etc) for producing high-quality schemes and graphics. If your graphical editor is Libre/Microsoft Office, export our chart as pdf/eps file.

Pay attention that the figure caption is placed \emph{below} the figure, while the table caption is placed \emph{above} the table (cf.~table~\ref{table}). Figures or tables may span both columns, if necessary (use \texttt{figure*} or \texttt{table*} environments).


\begin{table*}[tbp]
\caption{Defining characteristics of five early digital computers}\label{table}
\centering
\begin{tabular}{|@{\vrule width0ptheight9pt\enspace}l|c|c|c|c|c|c|c|}\hline
\hfil\bf Computer&\bf First operation&\bf Place& \bf Decimal/Binary&\bf Electronic&\bf Programmable&\bf Turing complete\\\hline
Zuse Z3&May 1941&Germany&binary&No&By punched film stock&Yes (1998)\\\hline
Atanasoff-Berry Computer&Summer 1941&USA&binary&Yes&No&No\\\hline
Colossus&1943--1944&UK&binary&Yes&Partially, by rewiring&No\\\hline
Harvard Mark I--IBM ASCC&1944&USA&decimal&No&By punched paper tape&Yes (1998)\\\hline
\multirow{2}{*}{ENIAC}&1944&USA&decimal&Yes&Partially, by rewiring&Yes\\\cline{2-7}
&1948&USA&decimal&Yes&By Function Table ROM&Yes\\\hline
\end{tabular}
\end{table*}

Position figures and tables at the tops and bottoms of columns.  Avoid placing them in the middle of columns.  Large figures and tables my span across both columns.  Figure captions should be below the figures; table captions should be above the tables.  Avoid placing figures and tables before their first mention in the text.  Use the abbreviation ``Fig.~1,'' even at the beginning of a sentence.


\section{Algorithms}

For algorithms typesetting the packages \verb|algorithmic| and~\verb|algorithm| are used. 

The first one defines environment \verb|algorithmic|, inside which algorithms are inputted. In the following example the code from the left side produces the output shown on the right side. The optional parameter~\verb|[2]| defines lines numeration of the algorithm.

The  package \verb|algorithm| defines new float environment \verb|algorithm| which is  similar to floats \verb|figure| and \verb|table|. 

The usage of the \verb|algorithm| environment is presented in the example~\ref{algorytm}, output---algorithm~\ref{example}---is respectively given on the page~\pageref{example}.

\begin{figure}
\begin{example}[the \texttt{algorithm} environment]\label{algorytm} output is given on the page~\pageref{example}.
\begin{minted}{latex}
\begin{algorithm}[tbp]
\caption{Calculating $x$ to the 
power~$n$\label{example}}
\begin{algorithmic}[2]
\REQUIRE $n\ge 0$
\ENSURE $a=x^n$
\STATE $k\leftarrow n$; $a\leftarrow 1$; 
\STATE $b\leftarrow x$;
\WHILE[Invariant: $x^n=a\cdot b^k$]
{$k>0$}
\IF{$k$ is even}
\STATE $k\leftarrow k/2$;
\STATE $b\leftarrow b\cdot b$;
\ELSE[$k$ is odd]
\STATE $k\leftarrow k-1$;
\STATE $a\leftarrow a\cdot b$;
\ENDIF
\ENDWHILE
\end{algorithmic}
\end{algorithm}
\end{minted}
\end{example}
\end{figure}

\begin{algorithm}[t]
\caption{Calculating $x$ to the 
power~$n$\label{example}}
\begin{algorithmic}[2]
\REQUIRE $n\ge 0$
\ENSURE $a=x^n$
\STATE $k\leftarrow n$; $a\leftarrow 1$; 
\STATE $b\leftarrow x$;
\WHILE[Invariant: $x^n=a\cdot b^k$]
{$k>0$}
\IF{$k$ is even}
\STATE $k\leftarrow k/2$;
\STATE $b\leftarrow b\cdot b$;
\ELSE[$k$ is odd]
\STATE $k\leftarrow k-1$;
\STATE $a\leftarrow a\cdot b$;
\ENDIF
\ENDWHILE
\end{algorithmic}
\end{algorithm}


\section{Code fragments}
To include the programming code fragment, please use the package \verb|minted|. This package ``understands'' almost all the programming languages and their dialects from BASIC to XML. Particulars of the package usage are given in the package documentation. Example~\ref{c++code} we shows implementation of the algorithm~\ref{example} in C++. Examples contain \TeX-code and the corresponding output.

\begin{figure}
\begin{example}[C++]\label{c++code}
output is shown on the page~\pageref{c++}.
\begin{minted}{latex}
\begin{minted}{c++}
\end{minted}
\begin{minted}{c++}
int power(int x,int n){
   int k,a,b;
   k=n; a=1; b=x;
   while(k>0){//Invariant: x^n=a*b^k
      if (k % 2==0){
         k/=2;
         b*=b;
      }
      else{
         k--;
         a*=b;
      }
   }
   return a;
}
\end{minted}
\mintinline{latex}|\end{minted}|
\end{example}
\end{figure}

\begin{algorithm}[tbp]
\caption{Calculating $x$ to the 
power~$n$ in C++\label{c++}}
\begin{minted}{c++}
int power(int x,int n){
   int k,a,b;
   k=n; a=1; b=x;
   while (k>0) {// Invariant: x^n=a*b^k
      if (k % 2==0){
         k/=2;
         b*=b;
      }
      else{
         k--;
         a*=b;
      }
   }
   return a;
}
\end{minted}
\end{algorithm}


\section{References}
For bibliography references use the standard \LaTeX\ tools. Refer simply to the reference number, as in~\cite{c3}. Do not use ``Ref.~\cite{c3}'' or ``reference~\cite{c3}'' except at the beginning of a sentence:  ``Reference~\cite{c3} was the first\dots''. Do not put footnotes in the reference list. IEEE Transactions no longer use a journal prefix before the volume number.  For example, use ``IEEE Trans.\ Magn., vol.~25,'' not ``vol.~MAG-25.''

Give all authors' names; do not use ``et al.'' unless there are six authors or more.  Papers that have not been published, even if they have been submitted for publication, should be cited as ``unpublished'' \cite{cu}.  Papers that have been accepted for publication should be cited as ``in press'' \cite{cin}.  Capitalize only the first word in a paper title, except for proper nouns and element symbols.
For papers published in translation journals, please give the English citation first, followed by the original language citation~\cite{ctr}.

\subsection{DOI}
Please, add a DOI information to your bibliography using the following format
\begin{minted}{latex}
\bibitem{citation}
Ghosh, M. K. and Harter, M. L.
\textit{A viral mechanism for remodeling 
  chromatin structure in G0 cells,} 
Mol. Cell. 12:255–260, 2003 
\url{https://dx.doi.org/10.1016/S1097}
\end{minted}


\subsubsection*{BibTeX users}
Since \verb|IEEEtran.bst| does not officially support doi yet, use the \verb|myIEEEtran.bst| by Gerald Q.~"Chip" Maguire Jr. (Available also at the conference site.) The template for doi in bib file is as follows:
\begin{minted}{bibtex}
@article{foo2010,
  author = "Mrinal K. Ghosh 
         and Marian L. Harter",
  journal = "Mol. Cell.",
  year = 2003,
  title = "A viral mechanism for 
   remodeling chromatin structure 
   in G0 cells",
  doi = {10.1016/S1097-2765(03)00225-9}
}
\end{minted}


\section{Some Common Mistakes}
 An excellent style manual and source of information for science writers is~\cite{c7}.

\section{Conclusion}
A conclusion section is not required. Although a conclusion may review the main points of the paper, do not replicate the abstract as the conclusion. A conclusion might elaborate on the importance of the work or suggest applications and extensions. 

\section{Submission}
Submitting the article, please, provide
\begin{itemize}
\item Reviewing:
\begin{itemize}
\item PDF file with the article. PDF will be sent to the referees.
\end{itemize}
\item Final version:
\begin{itemize}
\item PDF file with the article.
\item The source of the article.
\item Files of all the inserted graphics.
\item Bibtex files if any were used.
\end{itemize}
\end{itemize}

\section*{Appendix}
Appendixes, if needed, appear before the acknowledgment.

% serves to balance the column lengths on the last page of the document
% should be inserted the left column of the last page

% \balance

\section*{Acknowledgment}
The preferred spelling of the word ``acknowledgment'' in America is without an ``e'' after the ``g.''  Try to avoid the stilted expression, ``One of us (R.~B.~G.) thanks \dots'' Instead, try ``R.~B.~G.\ thanks\dots''  Put sponsor acknowledgments in the unnumbered footnote on the first page.

\begin{thebibliography}{9}


\bibitem{1}
Tian, Yahui, et al. "GDC-Net: a U-Net for precise brain vessel segmentation with global–local and depthwise attention plus content-aware upsampling." The Journal of Supercomputing 82.4 (2026): 215.

\bibitem{2}
Turjya, Sapthak Mohajon, and Mulham Fawakherji. "Federated lung nodule segmentation using a hybrid transformer–U-Net architecture." Scientific Reports (2026).

\bibitem{3}
Chauhan, Chhavi, et al. "The critical role of standards for AI in digital pathology: Digital Pathology Association Concept Paper." Journal of Pathology Informatics (2026): 100645.

\bibitem{4}
McMahan B, et al. Communication-Efficient Learning of Deep Networks from Decentralized Data. AISTATS, 2017.

\bibitem{5}
Bakas, Spyridon, et al. "Federated Learning in Healthcare: From Research to Real-World Deployment." Annual Review of Biomedical Engineering 28 (2026).

\bibitem{6}
Markodimitrakis, Emmanouil, et al. "Federated learning on magnetic resonance imaging: a critical review." Artificial Intelligence Review (2026)..

\bibitem{7}
Tanim, Sharia Arfin, et al. "FedFusionNet: Advancing Oral Cancer Recurrence Prediction through Federated Fusion Modeling." Information Fusion (2026): 104205..

\bibitem{LL1}
H. Guan, P. T. Yap, A. Bozoki, and M. Liu. "Federated learning for medical image analysis: A survey." Pattern recognition 151 (2024): 110424.

\bibitem{LL2}
S. Nazir, and M. Kaleem. "Federated learning for medical image analysis with deep neural networks." Diagnostics 13, no. 9 (2023): 1532.


\bibitem{LL3}
M. H. U. Rehman, W. H. L. Pinaya, P. Nachev, J. T. Teo, S. Ourselin, and M. J. Cardoso. "Federated learning for medical imaging radiology." The British Journal of Radiology 96, no. 1150 (2023): 20220890.


\bibitem{LL4}
D. Ghosh, M. Mehjabin, M. E. Rayed, M. F. Mridha, and M. M. Kabir. "Advancements and challenges of federated learning in medical imaging: a systematic literature review." Artificial Intelligence Review (2026).

\bibitem{LL5}
V. S. Parekh, S. Lai, V. Braverman, J. Leal, S. Rowe, J. J. Pillai, and M. A. Jacobs. "Cross-domain federated learning in medical imaging." arXiv preprint arXiv:2112.10001 (2021).

\bibitem{LL6}
S. Alphonse, F. Mathew, K. Dhanush, and V. Dinesh. "Federated learning with integrated attention multiscale model for brain tumor segmentation." Scientific Reports 15, no. 1 (2025): 11889.

\bibitem{LL7}
M. Sun, Z. Yang, Y. Huang, H. Yu, Y. Chen, S. Qi, A. B. J. Teoh, and Y. Zhang. "Federated learning for large models in medical imaging: a comprehensive review." arXiv preprint arXiv:2508.20414 (2025).

\bibitem{dataset}
L. Shi, X. Li, W. Hu, H. Chen, J. Chen, Z. Fan, M. Gao et al. "EBHI-Seg: a novel enteroscope biopsy histopathological hematoxylin and eosin image dataset for image segmentation tasks." Frontiers in Medicine 10 (2023): 1114673.

\bibitem{c1}
J. G. F. Francis, ``The QR Transformation I,'' {\it Comput.\ J.,} vol.~4, 1961, pp.~265--271.

\bibitem{c2}
H. Kwakernaak and R. Sivan, {\it Modern Signals and Systems,} Prentice Hall, Englewood Cliffs, NJ; 1991.

\bibitem{c3}
D. Boley and R. Maier, ``A Parallel QR Algorithm for the Non-Symmetric Eigenvalue Algorithm,'' {\it in Third SIAM Conference on Applied Linear Algebra,} Madison, WI, 1988, pp.~A20.


\bibitem{c7}
M. Young, {\it The Technical Writer's Handbook.} Mill Valley, CA:  University Science, 1989.

\bibitem{cg}
IEEE Guidelines for Author
Supplied Electronic
Text and Graphics, \url{http://www.ieee.org/portal/cms_docs_iportals/iportals/publications/journmag/transactions/eic-guide.pdf}


\bibitem{cu} K. Elissa, ``Title of paper if known,'' unpublished.
\bibitem{cin} ``R. Nicole, Title of paper with only first word capitalized,'' J.\ Name Stand.\ Abbrev., in press.
\bibitem{ctr} Y. Yorozu, M. Hirano, K. Oka, and Y. Tagawa, ``Electron spectroscopy studies on magneto-optical media and plastic substrate interface,'' IEEE Transl.\ J.\ Magn.\ Japan, vol.~2, pp.~740--741, August 1987 [Digests 9th Annual Conf. Magnetics Japan, p.~301, 1982].


\end{thebibliography}


\end{document}
